\documentclass[10pt,a4paper]{article}

\usepackage[a4paper,margin=1.05cm]{geometry}
\usepackage[hidelinks,unicode]{hyperref}
\usepackage{enumitem}
\usepackage{tabularx}
\usepackage{xcolor}
\usepackage{array}
\usepackage{graphicx}
\usepackage{fontspec}
\usepackage{polyglossia}

\setdefaultlanguage{russian}
\setotherlanguage{english}
\setmainfont{DejaVu Sans}
\definecolor{accent}{HTML}{1F5F8B}
\definecolor{textgray}{HTML}{4B5563}
\hypersetup{
  pdftitle={Александр Бородин -- ML Engineer | LLMs \& AI Agents},
  pdfauthor={Александр Бородин},
  pdfsubject={Резюме}
}

\pagestyle{empty}
\setlength{\parindent}{0pt}
\setlength{\parskip}{0pt}
\setlength{\tabcolsep}{3pt}
\emergencystretch=2em
\setlist[itemize]{leftmargin=1.0em,itemsep=0.06em,topsep=0.06em,parsep=0pt,partopsep=0pt}
\renewcommand\labelitemi{--}

\newcommand{\sectiontitle}[1]{%
  \vspace{0.38em}%
  {\normalsize\bfseries\color{accent} #1}\par
  \vspace{0.08em}\hrule\vspace{0.21em}%
}

\newcommand{\entry}[4]{%
  \textbf{#1}\hfill {\color{textgray}#2}\par
  {\color{textgray}#3}\par
  #4\vspace{0.21em}%
}

\newcommand{\achievement}[2]{%
  #1\hfill{\color{textgray}#2}\par
}

\begin{document}
\fontsize{9.00}{10.1}\selectfont
\raggedright

\begin{minipage}[c]{0.78\textwidth}
  {\LARGE\bfseries Александр Бородин}\par
  \vspace{0.04em}
  {\large ML Engineer | LLMs \& AI Agents}\par
  \vspace{0.16em}
  \href{mailto:apdorodin765@gmail.com}{apdorodin765@gmail.com} \quad
  +7 908 971 16 97 \quad
  \href{https://github.com/aleksanderborodin}{github.com/aleksanderborodin}\par
  \href{https://t.me/aleksanderbor}{Telegram: @aleksanderbor} \quad
  \href{https://aleksanderborodin.github.io/}{aleksanderborodin.github.io}
\end{minipage}
\hfill
\begin{minipage}[c]{0.15\textwidth}
  \raggedleft
  \setlength{\fboxsep}{0pt}
  \setlength{\fboxrule}{0.6pt}
  \fcolorbox{accent}{white}{\includegraphics[width=2.25cm]{assets/profile-original.jpg}}
\end{minipage}

\sectiontitle{Профиль}
ML-инженер, студент ФПМИ МФТИ. Разрабатываю системы на основе LLM: от RAG и интеграции моделей до мультиагентных систем и агентной эволюции — итеративного поиска и улучшения решений с помощью агентов. Есть опыт построения бенчмарков, оценки моделей и решения задач оптимизации.

\sectiontitle{Образование}
\entry{Бакалавриат, прикладная математика и физика}{2023--2027}{ФПМИ МФТИ, учусь; GPA 7.8/10}{Курсы: матанализ, теория вероятностей, линейная алгебра, алгоритмы, вычислительная сложность, ML, физика, экономика.}

\sectiontitle{Опыт работы}
\entry{ML-инженер — исследовательская практика}{февр.--авг. 2026 (6 мес.)}{ЦКМ МФТИ}{%
\begin{itemize}
  \item Разрабатывал агентов на основе LLM и бенчмарки для оценки качества их работы.
  \item Работал над агентной эволюцией: системами, в которых агенты итеративно генерируют, проверяют и улучшают решения.
  \item Участвовал в подготовке двух научных статей: одна находится на рецензировании в ACL, вторая готовится к подаче.
\end{itemize}
}

\entry{Product / AI Engineer}{янв.--май 2026}{\href{https://modelgate.ru}{ModelGate.ru} -- LLM API-шлюз и прокси-сервис}{%
\begin{itemize}
  \item Разрабатывал LLM API-шлюз с единым интерфейсом для доступа к моделям разных провайдеров.
  \item Интегрировал API провайдеров и реализовывал маршрутизацию запросов с выбором модели и провайдера.
  \item Создавал основные компоненты сервиса, включая backend и публичный сайт. Стек: Python, LLM API.
\end{itemize}
}

\entry{Разработчик моделей и алгоритмов}{сент. 2024--февр. 2025}{Розничный банк, клиентский проект}{%
\begin{itemize}
  \item Разработал алгоритм оптимизации инкассации; анализ исторических данных дал прогноз +20 млн руб. годовой прибыли.
  \item Построил пайплайн для скрапинга и анализа данных ЦБ РФ, недоступных через API, и внутреннюю RAG-систему для вопросов по документам.
  \item Стек: Python, NumPy, Pandas, RAG, скрапинг.
\end{itemize}
}

\sectiontitle{Проекты}
\entry{Idea Evolve}{2026}{\href{https://github.com/aleksanderborodin/idea_evolve}{github.com/aleksanderborodin/idea\_evolve}}{%
\begin{itemize}
  \item Эволюционный мультиагентный фреймворк: специализированные Claude-агенты работают параллельно по поколениям и строят общую базу знаний.
  \item Улучшил поиск множества Сидона с 66 до 105 элементов за 7 автономных поколений, выше ориентира $\sqrt{N}\approx100$ при $N=10000$.
  \item Реализовал ролевую специализацию агентов, возобновляемое состояние, извлечение знаний и Flask-дашборд.
  \item Стек: Python 3.12+, Claude Code CLI, OpenCode, Anthropic API, Flask, мультиагентная оркестрация.
\end{itemize}
}

\entry{Prefix Injection Jailbreak}{2025}{\href{https://github.com/aleksanderborodin/prefix_injection_jailbreak}{github.com/aleksanderborodin/prefix\_injection\_jailbreak}}{%
\begin{itemize}
  \item Исследовал префиксные атаки, вдохновлённые Choice Blindness: шаблоны заставляли Gemini 2.5 Pro и DeepSeek R1 завершать MaliciousInstruct-запросы почти в 100\% случаев против почти нулевого базового уровня.
  \item Стек: Python, LLM API, оценка моделей, проектирование промптов и контекста.
\end{itemize}
}

\entry{Yoola Explain}{2025}{\href{https://github.com/aleksanderborodin/yoola_explain}{github.com/aleksanderborodin/yoola\_explain}}{%
\begin{itemize}
  \item Браузерное расширение: извлекает Terms of Service и возвращает структурированное LLM-резюме — ключевые пункты, сбор данных и предупреждения.
  \item Добавил мультиязычность и кэширование. Стек: JavaScript, FastAPI, OpenRouter, LLM API.
\end{itemize}
}

\sectiontitle{Навыки}
\begin{tabularx}{\textwidth}{@{}>{\bfseries}p{2.50cm}X@{}}
Код & Python: продвинутый; C++: средний; чтение кода на C и Assembler. \\
Данные / ML & PyTorch: базовый; NumPy; Pandas; PostgreSQL; Matplotlib; ClearML; Excel. \\
NLP / LLM & BM25; эмбеддинги; RAG; LLM API; оценка моделей; проектирование промптов и контекста; тестирование джейлбрейков. \\
Агенты & Claude Code CLI; OpenCode; мультиагентная оркестрация; файловые процессы. \\
Бэкенд & FastAPI; Flask; LangChain; LangGraph; Linux: домашний сервер и рабочая станция. \\
Языки & Русский: родной; английский: B2+ по placement test МФТИ. \\
\end{tabularx}

\sectiontitle{Образовательные программы и достижения}
\entry{\href{https://airi.net/ru/summer-school-2026/}{«Лето с AIRI 2026» -- летняя школа по искусственному интеллекту}}{21 июля--4 авг. 2026}{Итоговый командный проект: \href{https://github.com/n4rly-boop/ideas-lake}{«Озеро идей»} -- долговременная память для эволюционных AI-систем с гибридным поиском и графом знаний.}{}
\achievement{\textbf{Золотая медаль}, \href{https://iepho.ru/}{Международная экспериментальная олимпиада по физике (IEPhO)}}{2023}
\achievement{\textbf{Призёр заключительного этапа} \href{https://vos.olimpiada.ru/}{ВсОШ по физике} — 10 и 11 классы}{2022, 2023}
\achievement{Победитель командных Wiki-гонок, Сириус}{2023}

\end{document}
